\section{Use Cases}\label{sec:use-cases}

Motivo suits composers who reason about repetition and variation but lack a concise editable source form.
The scenarios below are not exhaustive; they frame how the language and studio are intended to fit alongside
existing production tools.

\topic{Offline batch composition.}
A composer writes Motivo source, compiles once, and imports the resulting \texttt{.mid} into any host.
Deterministic output makes the artifact reproducible and easy to version alongside the source file.

\topic{Compact symbolic descriptions.}
Patterns, loops, and control flow can describe arrangements that would require hundreds of hand-placed events in a
piano roll.
The compiler expands those rules into a flat timeline at compile time, which is useful when repetitive structure
dominates the piece.

\topic{Teaching algorithmic composition.}
Motivo Studio pairs immediate diagnostics with piano-roll preview and browser playback, so students can see how a small
change in source affects the expanded score without configuring a synthesis stack.

\topic{Rapid prototyping before DAW polish.}
Motifs and harmonic sketches can be iterated in typed source and then refined manually in a digital audio workstation
for mixing, micro-timing, and sound design.
Motivo is meant to front-load structural work, not to replace downstream production steps.
